%%%PAGE 47 rot=0 legibility=good summary=带电粒子在偏转电场中的类平抛(加速+偏转+匀速直线，侧移公式D=y1+y2)；交变电压：正弦波(U临、有无射出时间比、屏上亮线长)与方形波(由E-t画v-t、终向相同/相反、垂直型、静止释放打板型)
%%%BLOCK id=047-01 ch=09 sec=偏转电场·类平抛 kind=derivation alt=none cont=none y=0.01-0.285
\red{①}\ \textbf{电场类平抛}

%FIG p047_a 0.25 0.005 0.558 0.195
\notefig[0.5]{figures/p047_a.png}
% 图内标注：U1、+、L1、d、−、θ、L2、y1、y2，及"匀加速""偏转""匀速直线运动"；图左侧另有粒子标注 m·q（在 U1 极板左边，未入图）
% 图右上方另有手写小注"板长．"及一个小箭头/L2 样的记号（未入图）

板长．

\[
\begin{cases}
v=at\ \to\ V=\sqrt{\dfrac{2qU_1}{m}}\\[6pt]
d=\dfrac{1}{2}at^{2}\ \to\ \sqrt{\dfrac{2d}{a}}=\sqrt{\dfrac{2md^{2}}{qU_1}}\\[6pt]
a=\dfrac{qU}{md}
\end{cases}
\]
% CHECK: 第二行根号内分母写作 qU1，按物理应为偏转电压 U2（或原稿即 U）

\[ qU_1=\frac{1}{2}mV_0^{2} \]
\[ y_1=\frac{1}{2}\,\frac{qU_2}{md}\left(\frac{L_1}{V_0}\right)^{2} \]
\[ \tan\theta=\frac{at}{V_0}=\frac{qU_2L_1}{mdV_0^{2}} \]
\[ \red{\star}\ D=y_1+y_2=\frac{1}{2}\,\frac{qU}{md}\left(\frac{L_1}{V_0}\right)^{2}+\frac{qU_2L_1L_2}{mdV_0^{2}}=\frac{U_2L_1^{2}}{4dU_1}+\frac{U_2L_1L_2}{2dU_1} \]
% CHECK: 第一项分子原稿为 qU（无下标，应为 qU_2）
\[ y_2=L_2\tan\theta \]

{\large 侧移与$U_{\text{偏}}$成正比}

侧移量与$K$无关（荷质比）故电\unclear{性}相同粒子同轨迹
%%%END
%%%BLOCK id=047-02 ch=09 sec=交变电场·正弦波 kind=method alt=none cont=none y=0.285-0.575
\red{②}\ \textbf{正弦波}（忽略偏转时间）

%FIG p047_b 0.05 0.325 0.24 0.435
\notefig[0.3]{figures/p047_b.png}
% 图内标注：U、U_m、U_0、t，以及横轴下方小字(t_有、t_无 一类)

\unclear{恰}好从极板边缘射出粒子的$U_{\text{临}}$

\[ \frac{d}{2}=\frac{U_0L_1^{2}}{4dU_1}\qquad \text{故}\ U_{\text{临}}=U_0=\frac{md^{2}v^{2}}{qL_1^{2}}=\frac{2d^{2}U_1}{L_1^{2}} \]

① 若 $U_m\le U_0$ 全飞出

② 若 $U_m>U_0$\quad $\dfrac{U_0}{U_m}=\dfrac{1}{2}$\ \ $\dfrac{t_{\text{有}}}{t_{\text{无}}}=\dfrac{1}{2}$\quad $\dfrac{U_{\text{\unclear{临}}}}{U_m}=\dfrac{\sqrt{2}}{2}$\ \ $\dfrac{t_{\text{有}}}{t_{\text{无}}}=\dfrac{1}{1}$\quad $\dfrac{U_{\text{\unclear{临}}}}{U_m}=\dfrac{\sqrt{3}}{2}$\ \ $\dfrac{t_{\text{有}}}{t_{\text{无}}}=\dfrac{2}{1}$
% CHECK: t_有/t_无 下标字形不甚清楚(按物理含义录为 有/无)；后两个分式分子写作 U 加一个汉字下标(似"临")，第一个为 U_0

\[
\text{屏上亮线长}\ \left\{\begin{array}{l@{\quad}l}
U_m\le U_0\text{时} & L=2D=2\Bigl(\dfrac{U_mL_1^{2}}{4dU_1}+\dfrac{U_mL_1L_2}{2dU_1}\Bigr)\\[10pt]
U_m>U_0\text{时} & \text{\unclear{必}有射出\ 有撞板}\ \ L=2D=2\Bigl(\dfrac{d}{2}+L_2\cdot 2\,\dfrac{\frac{d}{2}}{L_1}\Bigr)\\[10pt]
 & \qquad\qquad\ =d+L_2\tan\theta=d+2L_2\tan\theta
\end{array}\right.
\]
% CHECK: "d/2"分式的分母原稿写作 ℓ（疑为 L_1，下标可能漏写）；最末一个 tan 后的角度符号写得像 α/δ，此处按 θ 录入；"="后的 d 处有涂改液修改痕迹
%%%END
%%%BLOCK id=047-03 ch=09 sec=交变电场·方形波 kind=method alt=none cont=none y=0.575-0.845
\red{③}\ \textbf{方形波}\quad \fcolorbox{red!75!black}{white}{\begin{tabular}{c}$U$-$t$\\ 由 $E$-$t$ 画 $v$-$t$\end{tabular}}\quad 正上负下零化平．释放时刻画2周期

%FIG p047_c 0.04 0.604 0.415 0.853
\notefig[0.5]{figures/p047_c.png}
% 图内标注：E、V、t、T/2、T、3T/2、2T、L=nTV_0，红字"T/8 起步""3T/8 起步"（图为 E-t 方波与对应的 v-t 折线，红色虚线为另一释放时刻的 v-t）

%FIG p047_d 0.485 0.625 0.74 0.72
\notefig[0.4]{figures/p047_d.png}
% 图内红字：[0,T/4]段"终向与起步同向"；[T/4,3T/4]段"终向与反向起步"；[3T/4,T]段"终向与起步同向"

\[ \Bigl[0,\frac{T}{4}\Bigr]\cup\Bigl[\frac{3}{4}T,\,T\Bigr]\ \text{终向相同}\]
\[ \Bigl[\frac{T}{4},\frac{3}{4}T\Bigr]\ \text{终向相反．} \]

\[ \frac{T_1+T_2}{2}=\frac{T}{4}\ \ \text{终向相同} \]
\[ \frac{T_1+T_2}{2}=\frac{T}{2}\ \ \text{终向相反} \]
%%%END
%%%BLOCK id=047-04 ch=09 sec=交变电场·方形波 kind=method alt=none cont=none y=0.845-1.0
① \textbf{垂直型}

%FIG p047_e 0.02 0.848 0.29 0.896
\notefig[0.4]{figures/p047_e.png}
% 图内标注：nT、V_0（多个向右箭头）

周期整数倍的运动时间\ 出射速度与入射速度相同．

② \ 静止释放打板型．

%FIG p047_f 0.05 0.935 0.20 0.995
\notefig[0.25]{figures/p047_f.png}
%%%END
%%%PAGEEND 47
